\subsection{拓扑群的庞加莱群}

若 G 是拓扑群，则对每个 \(g \in G\)，左平移 \(x \mapsto gx\) 都是 G 到自身的同胚（TG, III, 第~2~页），并诱导出从 \(\pi_{1}(G, e)\) 到 \(\pi_{1}(G, g)\) 的同构。

设 G 是拓扑群，e 是其单位元，并以 \(G_{0}\) 表示 G 中 e 的道路连通分支。以 R 表示 G 上其等价类为 G 的道路连通分支的等价关系，并令 \(p: G \to \pi_{0}(G)\) 为规范映射。由于每个左平移或右平移都是 G 的同胚，关系 R 对 G 的群法则在左、右两侧都相容。根据 A, I, 第~35~页的定理 2，群 \(G_{0}\) 是 G 的正规子群，商群 \(G/G_{0}\) 等于赋予 G 所诱导的合成法则的 \(\pi_{0}(G)\)，且映射 p 是群同态。

群 \(\pi_{1}(G,e)\) 称为 G 的庞加莱群，记作 \(\pi_{1}(G)\)。从 \(G_{0}\) 到 G 的规范嵌入诱导出从 \(\pi_{1}(G_{0})\) 到 \(\pi_{1}(G)\) 的同构（III, 第~293~页，注 2）。

设 \(g\) 是 G 的一个元素；右平移 \(\boldsymbol{\delta}_g \colon x \mapsto xg\) 诱导出同构 \((\boldsymbol{\delta}_g)_* \colon \pi_1(\mathrm{G}) \to \pi_1(\mathrm{G}, g)\)。设 \(g\) 属于 \(\mathrm{G}_0\)，并令 \(b\) 为连接 \(e\) 与 \(g\) 的道路。映射 \(\sigma \colon \mathrm{G} \times \mathbf{I} \to \mathrm{G}\) 由 \(\sigma(x, t) = xb(t)\) 定义，是将 \(\mathrm{Id}_{\mathrm{G}}\) 与 \(\boldsymbol{\delta}_g\) 连接起来的同伦。根据 III, 第~295~页的命题 3，对于每个 \(\alpha \in \pi_1(\mathrm{G})\)，有 \((\boldsymbol{\delta}_g)_*(\alpha) = \beta^{-1}\alpha\beta\)，其中 \(\beta \in \varpi_{e,g}(\mathrm{G})\) 是道路 \(b\) 的类。若以 \(\gamma_g\) 表示左平移 \(x \mapsto gx\)，同样有 \((\gamma_g)_*(\alpha) = \beta^{-1}\alpha\beta\)。

对每个 \(g \in G\)，映射 \(\operatorname{Int}(g)\) : \(G \to G\) 都诱导出 \(\pi_{1}(\operatorname{Int}(g))\) 在 \(\pi_{1}(G)\) 上的一个自同构。于是定义 G 在 \(\pi_{1}(G)\) 上的作用。对于所有 \(g \in G\)，有 \(\pi_{1}(\operatorname{Int}(g)) = (\boldsymbol{\gamma}_{g^{-1}})_{*} \circ (\boldsymbol{\delta}_{g})_{*}\)，因此子群 \(G_{0}\) 的作用是平凡的；这就给出了 \(\pi_{0}(G)\) 在 \(\pi_{1}(G)\) 上的作用。当 G 是交换群时，此作用是平凡的，但并非总是如此（IV, 第~459~页，习题 1）。

